\section{Problem 7}
\label{sec:problem7}

A spacecraft orbits the Earth, which is treated as an inertial reference frame, with orbital elements
\[
a = 15307.548\ \text{km}, \quad e = 0.7, \quad \Omega = 194^\circ, \quad i = 39^\circ, \quad \omega = 85^\circ, \quad \nu = 48^\circ,
\]
and $\mu = 398600$ km$^3\cdot$s$^{-2}$. The position and inertial velocity in ECI coordinates are obtained with the Problem 2 function. The orbit is shown in \cref{fig:p7-10}.

\programoutput[firstline=3, lastline=4]{outputs/problems7to10.txt}

\noindent The periapsis radius, $r_p = a(1-e) \approx 4592$ km, is smaller than $R_E = 6378$ km, so this orbit intersects the Earth. The given position, $r \approx 5317$ km, already lies inside the Earth.

\FloatBarrier
